\documentclass[10pt]{article}
\usepackage[T1]{fontenc}
\usepackage[utf8]{inputenc}
\usepackage[french]{babel}
\usepackage{amsmath,amssymb,mathrsfs,mathpazo}
\usepackage{geometry}
\geometry{a4paper,margin=2cm}
\usepackage{graphicx,array,tabularx,tabularray,colortbl,multirow,multicol,booktabs,enumitem,pifont,ulem,fancybox,diagbox,cancel}
\usepackage{pstricks,pstricks-add,pst-plot,pst-tree,pst-node,pst-text,pst-func,pst-math,pst-bspline,pst-3dplot,pst-eucl,pst-solides3d}
\let\pscalculate\undefined
\usepackage{pst-bezier}
\usepackage{tikz,pgfplots,tkz-tab}
\usetikzlibrary{arrows,arrows.meta,patterns,calc,positioning,trees,decorations.markings,shapes.geometric,shapes.misc,shapes.arrows,intersections,decorations.pathmorphing,decorations.pathreplacing,backgrounds,shadows,babel}
\pgfplotsset{compat=1.18}
\usepgfplotslibrary{fillbetween}
\usepackage[np]{numprint}
\usepackage{listings,makecell,pas-tableur,twoopt}
\let\barre\undefined
\usepackage[pstricks]{tablvar}
\definecolor{darkgreen}{RGB}{0,100,0}
\providecommand{\red}{\color{red}}
\providecommand{\blue}{\color{blue}}
\providecommand{\green}{\color{green}}
\providecommand{\black}{\color{black}}
\providecommand{\white}{\color{white}}
\providecommand{\gray}{\color{gray}}
\providecommand{\cyan}{\color{cyan}}
\providecommand{\magenta}{\color{magenta}}
\providecommand{\yellow}{\color{yellow}}
\providecommand{\violet}{\color{violet}}
\providecommand{\orange}{\color{orange}}
\providecommand{\euro}{\texteuro}
\providecommand{\EUR}{\texteuro}
\providecommand{\ieme}{\textsuperscript{e}}
\providecommand{\degres}{\ensuremath{^\circ}}
\providecommand{\degree}{\ensuremath{^\circ}}
\providecommand{\vv}[1]{\overrightarrow{#1}}
\providecommand{\Coord}[1]{#1}
\providecommand{\decosix}{}
\providecommand{\starredbullet}{\ensuremath{\bullet}}
\pagestyle{empty}
\setlength{\parindent}{0pt}
\usepackage{adjustbox}\DeclareUnicodeCharacter{2236}{:}\DeclareUnicodeCharacter{2190}{\ensuremath{\leftarrow}}
\providecommand{\R}{\mathbb{R}}
\providecommand{\N}{\mathbb{N}}
\providecommand{\Z}{\mathbb{Z}}
\providecommand{\Q}{\mathbb{Q}}
\providecommand{\D}{\mathbb{D}}
\providecommand{\C}{\mathbb{C}}
\providecommand{\e}{{\mathrm{e}}}
\providecommand{\vect}[1]{\overrightarrow{#1}}
\providecommand{\vectt}[1]{\overrightarrow{\mathrm{#1}}}
\providecommand{\vvt}[1]{\overrightarrow{\mathrm{#1}}}
\providecommand{\np}[2][]{{#2~#1}}
\providecommand{\vv}[1]{{\overrightarrow{#1}}}
\renewcommand{\d}{{\mathrm{d}}}
\providecommand{\degres}{{^\circ}}
\providecommand{\degre}{{^\circ}}
\providecommand{\blue}{\color{blue}}
\providecommand{\red}{\color{red}}
\providecommand{\enskip}{\quad}
\providecommand{\llbracket}{\lbrack\!\lbrack}
\providecommand{\rrbracket}{\rbrack\!\rbrack}
\providecommand{\euro}{{\text{€}}}
\providecommand{\no}{{\text{n°}}}
\providecommand{\interval}[2]{{[#1;#2]}}
\providecommand{\intervFO}[2]{{[#1;#2[}}
\providecommand{\dotfill}{\ldots}

\providecommand{\ds}{\displaystyle}
\providecommand{\textding}[1]{#1}
\providecommand{\Oij}{\(\mathrm{O};\vec{\imath},\vec{\jmath}\)}
\providecommand{\Oijk}{\(\mathrm{O};\vec{\imath},\vec{\jmath},\vec{k}\)}
\providecommand{\Ouv}{\(\mathrm{O};\vec{u},\vec{v}\)}
\providecommand{\intervFF}[2]{\ensuremath{[#1;#2]}}
\renewcommand{\euro}{\ifmmode\text{\texteuro}\else\texteuro\fi}
\DeclareUnicodeCharacter{00B0}{\ifmmode^{\circ}\else\textdegree\fi}
\DeclareUnicodeCharacter{20AC}{\ifmmode\text{\texteuro}\else\texteuro\fi}
% Énoncés : Éducation nationale. Transcription et corrigés : APMEP, Denis Vergès et contributeurs. Découpage et classement : CoopMaths, Sébastien Lozano, Éric Elter et contributeurs. Adaptation de la présentation : Maths Explorer.
% CC BY-SA : https://coopmaths.fr/www/about/
% https://coopmaths.fr/alea/static/e3c/2020/tex/e3c_2020_00_sujet35_1.tex
\begin{document}
\noindent\textit{APMEP / CoopMaths. CC BY-SA. Présentation adaptée par Maths Explorer.}\par\medskip
\medskip

Cet exercice est un questionnaire à choix multiples. Pour chaque question une seule réponse est exacte. Une mauvaise réponse ou une absence de réponse n'enlève aucun point. La bonne réponse rapporte un point. Il n'est pas demandé de justification.

\begin{enumerate}
\item  L'ensemble des solutions de l'inéquation $3x^2-4x+1\geqslant 0$ est :

\medskip

\begin{tabularx}{\linewidth}{*{2}{X}}
\textbf{A.~~} $]-\infty~;~-1]\cup[-\frac{1}{3}~;~+\infty[ $ &\textbf{B.~~} $ 
\left]-\infty~;~\frac{1}{3}\right]\cup [1~;~+\infty[$\\\textbf{C.~~}$\left]-\infty~;~-\frac{1}{3}\right]\cup\left[1~;~+\infty\right[ $& \textbf{D.~~} $ \left[\frac{1}{3}~;~1\right] $.
\end{tabularx}

\item Dans le plan muni d'un repère orthonormé, on considère les vecteurs $\vv{u} \ \binom{a+2}{-1}$ et $\vv{v}\ \binom{3}{a}$, où $a$ est un nombre réel. Les vecteurs $\vv{u}$ et $\vv{v}$ sont orthogonaux si, et seulement si :

\medskip

\begin{tabularx}{\linewidth}{*{4}{X}}
\textbf{A.~~} $ a(a+2)-3=0$ &\textbf{B.~~} $a(a+2)+3=0 $&\textbf{C.~~}$3(a+2)-a=0 $& \textbf{D.~~} $3(a+2)+a=0  $.
\end{tabularx}

\item Dans le plan muni d'un repère orthonormé, on considère le point A $(-2~;~3)$ et le vecteur $\vv{u} (1~;~2)$. Une équation cartésienne de la droite $d$ passant par le point A et de vecteur normal $\vv{u}$ est :

\medskip

\begin{tabularx}{\linewidth}{*{4}{X}}
\textbf{A.~~} $-2x+y-7=0 $ &\textbf{B.~~} $x+2y-4=0 $&\textbf{C.~~}$x-2y+8=0 $& \textbf{D.~~} $2x+y+1=0  $.
\end{tabularx}

\item On considère la suite $\left(u_n\right)$, géométrique de raison 2 et de premier terme $u_0=3$.

La somme $u_0+u_1+\dots+u_{10}$ est égale à :

\medskip

\begin{tabularx}{\linewidth}{*{4}{X}}
\textbf{A.~~} $3\left(2^{11}-1\right) $ &\textbf{B.~~} $3\left(1-2^{11}\right) $&\textbf{C.~~}$3\left(2^{10}-1\right) $& \textbf{D.~~} $ 3\left(1-2^{10}\right) $.
\end{tabularx}


\item Soit $f$ la fonction définie et dérivable sur $]1~;~+\infty[$ par $f(x)=\dfrac{2x+1}{x-1}$.

La fonction dérivée de $f$ sur $]1~;~+\infty[$ a pour expression :

\medskip

\begin{tabularx}{\linewidth}{*{4}{X}}
\textbf{A.~~} {\small $f'(x)=-\dfrac{1}{(x-1)^2}$} &\textbf{B.~~} {\small $f'(x)=-\dfrac{3}{(x-1)^2} $}&\textbf{C.~~}{\small $f'(x)=\dfrac{4x-1}{(x-1)^2}$}& \textbf{D.~~} {\small $f'(x)=\dfrac{1}{(x-1)^2}$}.
\end{tabularx}
\end{enumerate}

\vspace{0,5cm}

\end{document}
